% Part: first-order-logic
% Chapter: tableaux
% Section: derivations

\documentclass[../../../include/open-logic-section]{subfiles}

\begin{document}

\iftag{FOL}
      {\olfileid[pt]{fol}{tab}{der}}
      {\olfileid[pt]{pl}{tab}{der}}

\olsection{\usetoken{P}{tableau}}

\begin{explain}
Já dissemos o que é uma hipótese e apresentamos as regras de
inferência.  !!^{tableau}s são gerados indutivamente a partir delas:
cada !!{tableau} ou é um único ramo composto por uma ou mais
hipóteses, ou resulta de !!a{tableau} pela aplicação de uma das regras
de inferência a um ramo.
\end{explain}

\begin{defn}[!!^{tableau}]
!!^a{tableau} para as hipóteses $\sFmla{S_1}{!A_1}$, \dots,
$\sFmla{S_n}{!A_n}$ (em que cada $S_i$ é $\True$ ou~$\False$) é
uma árvore finita de !!{signed formula}s que satisfaz as seguintes condições:
\begin{enumerate}
\item As $n$ !!{signed formula}s no topo da árvore são
  $\sFmla{S_i}{!A_i}$, uma abaixo da outra.
\item Toda !!{signed formula} na árvore que não seja uma das
  hipóteses resulta da aplicação correta de uma regra de inferência
  a !!a{signed formula} no ramo acima dela.
\end{enumerate}
Um ramo de !!a{tableau} é \emph{fechado} sse contém tanto
$\sFmla{\True}{!A}$ quanto~$\sFmla{\False}{!A}$, e é \emph{aberto}
caso contrário. !!^a{tableau} em que todo ramo é fechado é um
\emph{!!{tableau} fechado} (para seu conjunto de hipóteses). Se !!a{tableau} não
é fechado, isto é, se contém pelo menos um ramo aberto, ele é
\emph{aberto}.
\end{defn}

\begin{ex}
  Todo conjunto de hipóteses, por si só, é !!a{tableau}, mas em geral
  não será fechado. (Evidentemente, ele só é fechado se as hipóteses
  já contêm um par de !!{signed formula}s
  $\sFmla{\True}{!A}$ e~$\sFmla{\False}{!A}$.)

  A partir de !!a{tableau} (aberto ou fechado), podemos obter outro,
  novo e maior, aplicando uma das regras de inferência a
  !!a{signed formula}~$!A$ nele. A regra acrescentará uma ou mais
  !!{signed formula}s ao final de todo ramo que contenha a ocorrência
  de~$!A$ à qual aplicamos a regra.

  Por exemplo, considere a hipótese $\sFmla{\True}{!A \land \lnot
    !A}$. Este é o !!{tableau} (aberto) composto apenas por essa
  hipótese:
  \begin{center}
    \begin{tableau}{}
      [\sFmla{\True}{\formula{A} \land \lnot \formula{A}}, just=\TAss]
    \end{tableau}{}
  \end{center}
  Obtemos dele um novo !!{tableau} aplicando a regra $\TRule{\True}{\land}$
  à hipótese. Essa regra permite acrescentar duas novas linhas ao
  !!{tableau}, $\sFmla{\True}{!A}$ e $\sFmla{\True}{\lnot !A}$:
  \begin{center}
    \begin{tableau}{}
      [\sFmla{\True}{\formula{A} \land \lnot \formula{A}}, just=\TAss
        [\sFmla{\True}{\formula{A}}, just={\TRule{\True}{\land}[1]},
          [\sFmla{\True}{\lnot \formula{A}}, just={\TRule{\True}{\land}[1]}
          ]
        ]
      ]
    \end{tableau}{}
  \end{center}
  Ao escrever !!{tableau}s, registramos à direita as regras que
  aplicamos (por exemplo, $\TRule{\True}{\land} 1$ significa que a
  !!{signed formula} nessa linha resulta da aplicação da regra
  $\TRule{\True}{\land}$ à !!{signed formula} da linha~$1$).
  Esse novo !!{tableau} contém agora !!{signed formula}s adicionais,
  mas só podemos aplicar uma regra a uma delas
  ($\sFmla{\True}{\lnot !A}$), neste caso, a regra
  $\TRule{\True}{\lnot}$. Isso resulta no !!{tableau} fechado
  \begin{center}
    \begin{tableau}{}
      [\sFmla{\True}{\formula{A} \land \lnot \formula{A}}, just=\TAss
        [\sFmla{\True}{\formula{A}}, just={\TRule{\True}{\land}[1]}
          [\sFmla{\True}{\lnot \formula{A}}, just={\TRule{\True}{\land}[1]}
            [\sFmla{\False}{\formula{A}}, just={\TRule{\True}{\lnot}[3]}, close]
          ]
        ]
      ]
    \end{tableau}{}
  \end{center}
\end{ex}

\end{document}
